%% notes-tex/figure-3-gate/sections/7-predictions.tex
%% SOURCE: every number here is in experiments/019-simb-multimodal/results/prediction_scatter.json,
%% written by experiments/019-simb-multimodal/scripts/prediction_scatter.py from the v13
%% validation prediction dump ($DATA_ROOT/val-predictions, run wq8y8nd5) and the fig3_core
%% split-seed-0 records; the figure is the script's SVG converted by `make plots`.

\clearpage
\section{What the expression predictions look like\secstatus{todo}}

The scores above are per-gene Pearson summaries. Figure~\ref{fig:predscatter} shows the
predictions themselves for the best expression checkpoint on file, the v13 reference on
split seed 0 (validation Pearson per gene 0.223), beside the bilinear ridge (0.134) and the
nearest-neighbor average (0.130) on ProtT5, refit on the same 155 validation strains and
6{,}127 graph genes with the cells selected on validation.

\pfstep{The pooled correlation is mostly the gene means} The 0.45 of panel a is not a model
score. The per-gene train mean alone, repeated for every strain, correlates 0.40 with the
measurements across all pairs, because genes differ in their mean log2 ratio (the gene means
have sd 0.09 against a pooled sd of 0.22 and carry 16.7 percent of the pooled variance), and
the model's predictions correlate 0.77 with that mean profile. With the gene means removed
from both sides (panel b) the pooled correlation is 0.28 and the mean per-gene Pearson is
0.22.

\pfstep{Direction, not distance} Within a gene the model's deviations have a third of the
measured spread (median sd ratio 0.32, panel e) and a regression slope of 0.09 on the
measured value. The nearest-neighbor average has the full spread (1.08) at a lower
correlation (0.17). The ratio that minimizes squared error is the Pearson itself, 0.22, so
the model is 1.4 times too spread for squared error and three times too narrow against the
measurements; both statements describe the same 0.32. Per-gene Pearson rises with how much a
gene moves (panel d): from 0.1 on the quietest genes to about 0.4 on the most variable for
the model, 0.1 to 0.2 for both baselines.

\pfstep{No predictor reaches the large responses} The strongest validation strain (panel f,
the YER068W deletion) spans measured log2 ratios from $-2$ to $+2.5$ across genes; neither
predictor leaves $\pm 0.5$. A median strain (panel g) is the gene-mean profile with a small
offset. The model's best gene (panel h, YKL145W, Pearson 0.62) has its strains in the right
order over a predicted range of 0.0 to 0.15 against a measured $-1$ to $0.5$.

\tcfig{figures/prediction_scatter_2026-10-06-17-47-55.pdf}{\textbf{The v13 reference
predicts the direction a gene moves across strains and not the distance; the large responses
are missing from every predictor.} Validation strains of split seed 0 on fig3\_core (155
strains, 6{,}127 graph genes). CGT: v13 reference checkpoint, run wq8y8nd5; ridge and
neighbor average: B2 and B3 on ProtT5, cells selected on validation.
\textbf{a}, Predicted against measured log2 ratio, every strain-gene pair (hexbin, log
count); pooled Pearson 0.45. \textbf{b}, The same with each gene's train mean removed from
both axes; pooled Pearson 0.28, regression slope 0.09 (orange). \textbf{c}, The neighbor
average on the axes of b; 0.17, slope 0.18. \textbf{d}, Per-gene Pearson over the 155
strains against the gene's train sd, with the running median per predictor (40 bins);
means 0.223, 0.134, 0.130. \textbf{e}, Per-gene sd(predicted) / sd(measured) on the same
axis; the solid line is 1, the dotted line 0.22 is the model's mean Pearson; medians 0.32,
0.56, 1.08. \textbf{f}, \textbf{g}, One strain each, genes ordered by the measured value:
the strongest-response strain (YER068W deletion, sd over genes 0.56, per-strain Pearson
0.26) and a median one (YJL187C deletion, 0.17, 0.17). \textbf{h}, \textbf{i}, One gene
each over the 155 strains: the gene with the highest model Pearson (YKL145W, 0.62; neighbor
average 0.58) and the gene at the model's median (YHL028W, 0.22; 0.06). Dashed lines are
the identity. One checkpoint, one partition; no uncertainty is shown.}{fig:predscatter}

\begin{tabular}{@{}p{8mm}p{170mm}@{}}
%% SOURCE: experiments/019-simb-multimodal/results/prediction_scatter.json (prediction_scatter.py)
\toprule
panel & files and keys \\
\midrule
all & \file{prediction_scatter.json}: \texttt{predictors} (per-gene mean and median Pearson, pooled raw and centered Pearson, median spread ratio, centered slope), \texttt{pooled\_decomposition}, \texttt{examples}; predictions from \file{val-predictions/compute-0-2-2397312\_3fcfb279\ldots.json}; baseline cells from \file{expression\_baselines\_split/seed0.json} \\
\bottomrule
\end{tabular}
